\section{OrchLang}\label{sec:language}

OrchLang is a small, statically typed language for LLM workflows. A workflow
declares its inputs, secrets, models, prompt templates and tools, then calls
models, branches, retries, invokes tools and returns one output (\cref{fig:syntax}).
The compiler, \code{orchc}, makes no network requests; models are named by
provider and model name and are only ever called by the mock runtime described in
\cref{sec:implementation}.

\begin{figure}[t]
\centering\footnotesize
\begin{tabular}{@{}r@{\;}c@{\;}l@{}}
$w$ & ::= & \code{workflow} $X$ \code{budget} $n$ [\code{leaks} $b$] [\code{observer} $O$] \code{\{} $d^*\ s^*$ \code{\}} \\[2pt]
$d$ & ::= & \code{input} $x$\code{:} $T$ [\code{untrusted}] $\ell^*$\code{;}
      $\mid$ \code{secret} $x$\code{:} $T$ $\ell^*$\code{;} \\
    & $\mid$ & \code{model} $x$ \code{= mock(}"\textit{provider}\code{/}\textit{name}"\code{) max\_tokens} $n$ \\
    &        & \quad [\code{tokenizer} $\tau$] [\code{overhead} $n$] [\code{max\_bytes} $n$]\code{;} \\
    & $\mid$ & \code{prompt} $p$\code{(}$x_1$\code{:} $T_1$, \dots\code{) -> } $T$ \code{=} "\textit{template with} \code{\{}$x_i$\code{\}}"\code{;} \\
    & $\mid$ & \code{tool} $f$\code{(}$x_1$\code{:} $T_1$, \dots\code{);} \\
$\ell$ & ::= & \code{max\_bytes} $n$ $\mid$ \code{max\_tokens} $n$ \code{tokenizer} $\tau$ $\mid$ \code{max\_tokens} $n$ \\[2pt]
$s$ & ::= & \code{let} $y$\code{:} $T$ \code{= call} $p$\code{(}$a_1$, \dots, $a_k$\code{) using} $x$\code{;} \\
    & $\mid$ & \code{emit} $f$\code{(}$a_1$, \dots, $a_k$\code{);}
      $\;\mid\;$ \code{output} $a$\code{;} \\
    & $\mid$ & \code{if} $g$ \code{\{} $s^*$ \code{\}} [\code{else \{} $s^*$ \code{\}}]
      $\;\mid\;$ \code{retry} $n$ \code{\{} $s^*$ \code{\}} \\
    & $\mid$ & (\code{declassify} $\mid$ \code{endorse})\code{(}$x$\code{) as} $y$\code{:} $T$ \code{because} "\textit{justification}"\code{;} \\
$g$ & ::= & $x$ $\mid$ \code{tokens(}$x$\code{)} $\bowtie$ $n$ \qquad ($\bowtie\ \in \{\code{<}, \code{<=}, \code{>}, \code{>=}, \code{==}, \code{!=}\}$) \\
$a$ & ::= & $x$ $\mid$ "\textit{literal}" \\
\end{tabular}
\caption{Abridged syntax of OrchLang. The full grammar, including \code{require} statements and
monetary prices, is in the language specification that accompanies the artifact.}
\label{fig:syntax}
\end{figure}

\subsection{Design commitments}

Three decisions make the analyses decidable, and \cref{sec:eval-real} measures what
they cost. First, the shape of a workflow is fixed before it runs: the only loop is
\code{retry} $n$ with a literal $n$, and the only branch tests a declared value.
Second, every length the compiler cannot see is declared in a stated unit: bytes
(\code{max\_bytes}), tokens of a named tokenizer (\code{max\_tokens} $n$
\code{tokenizer} $\tau$), or, for an estimate only, tokens. Third, leaving the
security lattice requires a written justification (\code{declassify},
\code{endorse}), which the compiler copies into its certificate.

Every declaration receives a unique binding identity, and all analyses refer to
bindings, models and prompts by identity, never by name. A redeclaration inside a
branch arm is therefore a different binding. This choice was forced by one of the
counterexamples in \cref{sec:failed}. Inputs and secrets may be declared only at
the top level of a workflow, since an input declared inside an arm would have no
single meaning for the environment that supplies it.

\subsection{Labels}\label{sec:labels}

Values carry labels from the product lattice
$\{\Pub \sqsubseteq \Sec\} \times \{\Tru \sqsubseteq \Unt\}$ of confidentiality
(public, secret) and integrity (trusted, untrusted)~\cite{denning1976lattice}; we
write $\conf(\ell)$ and $\integ(\ell)$ for the two components and $\sqcup$ for the
join. A literal and an ordinary input are $(\Pub,\Tru)$, an input declared
\code{untrusted} is $(\Pub,\Unt)$, and a secret is $(\Sec,\Tru)$. A model's
response is labelled with the join of everything that reached its request, so text
from a retrieved page stays untrusted through any number of calls. It can drive an
effect only after someone endorses it and writes down why.

Each binding also carries a \emph{data label}, which records what its content
depends on, separately from the program counter $pc$ under which it was computed.
The distinction matters inside a secret-guarded arm. A program-counter discipline
labels every value computed there as secret, so the second call of a chain inside
the arm could not receive the first call's response. Yet the content of that
response depends only on public data; only its existence depends on the secret,
and existence is exactly what the relational analysis of \cref{sec:requests}
accounts for. OrchLang therefore checks prompt arguments against the data label,
while outputs and effects, whose occurrence is visible outside the request trace,
are still checked against the program counter. \Cref{fig:rules} gives the rules for
calls, branches, effects and outputs.

\begin{figure}[t]
\centering\footnotesize
\begin{mathpar}
\inferrule*[right=Call]
  {\Gamma(a_i) = (\ell_i, \delta_i) \\ \conf(\delta_i) = \Pub \ \ (1 \le i \le k) \quad\text{[E230]}}
  {\Gamma; pc \vdash \code{let}\ y = \code{call}\ p(a_1,\dots,a_k)\ \code{using}\ m \\
   \dashv \Gamma[y \mapsto (pc \sqcup \textstyle\bigsqcup_i \ell_i,\;
   (\Pub, \integ(pc)) \sqcup \textstyle\bigsqcup_i \delta_i)]}
\and
\inferrule*[right=If]
  {\Gamma(x) = (\ell, \delta) \\ \Gamma; pc \sqcup \ell \vdash s_1 \\ \Gamma; pc \sqcup \ell \vdash s_2}
  {\Gamma; pc \vdash \code{if}\ x\ \{s_1\}\ \code{else}\ \{s_2\} \dashv \Gamma}
\and
\inferrule*[right=Emit]
  {\Gamma(a_i) = (\ell_i, \delta_i) \\ \ell_i \sqsubseteq (\Pub, \Tru) \quad\text{[E232, E233]} \\
   pc \sqsubseteq (\Pub, \Tru) \quad\text{[E234, E235]}}
  {\Gamma; pc \vdash \code{emit}\ f(a_1,\dots,a_k) \dashv \Gamma}
\and
\inferrule*[right=Output]
  {\Gamma(a) = (\ell, \delta) \\ \conf(\ell \sqcup pc) = \Pub \quad\text{[E231]}}
  {\Gamma; pc \vdash \code{output}\ a \dashv \Gamma}
\end{mathpar}
\caption{Selected label rules. $\Gamma$ maps each binding to its label $\ell$ and data label
$\delta$; literals are $(\Pub,\Tru)$. A premise that fails produces the diagnostic in brackets.
In \textsc{Call} the data label ignores the confidentiality of $pc$ but keeps its integrity. In
\textsc{If}, a guard whose data label is secret also creates the relational obligation of
\cref{sec:requests}.}
\label{fig:rules}
\end{figure}

The data label keeps integrity joined with the program counter, which is
conservative: a value computed inside an arm guarded by untrusted data is untrusted
even if it was computed only from trusted data. The real-workflow evaluation found
two cases where this over-reports (\cref{sec:eval-real}). Reclassification copies a
value into a new binding and relabels one component: \code{declassify} lowers
confidentiality to that of the program counter and \code{endorse} lowers integrity
to that of the program counter, so relabelling inside a secret arm still yields a
secret-dependent value. An endorsed copy of a secret keeps a secret data label.
Because the only bindings with secret data labels are secrets and their endorsed
copies, and \textsc{Call} keeps both out of requests, secret content never appears
in a request.

\begin{example}\label{ex:triage-labels}
In \cref{fig:triage}, \code{reply} is $(\Pub,\Unt)$ because the untrusted
\code{ticket} reached its request. Inside the arm, $pc = (\Sec,\Tru)$, so
\code{second} gets label $(\Sec,\Unt)$ but data label $(\Pub,\Unt)$: another
call could still receive it. Without the \code{endorse} on lines 16--17,
\textsc{Emit} fails with \code{E233}, because text derived from the ticket would
choose what the tool sends. With it, \code{vetted} is $(\Pub,\Tru)$ and the flow
check passes; the branch on \code{high\_risk} is then the relational analysis's
concern.
\end{example}

A workflow that passes every check receives a certificate: a JSON document bound to
the source by its SHA-256 digest, which records the token bounds and their
derivation, every label, every reclassification with its justification, and the
leakage report. The certificate is a report, not a proof object; no independent
checker exists yet (\cref{sec:conclusion}).
